\chapter*{คำนำ}
\addcontentsline{toc}{chapter}{คำนำ}

ศตวรรษที่ 20 ถือเป็นยุคทองแห่งการปฏิวัติทางฟิสิกส์ เมื่อมโนทัศน์ดั้งเดิมของกลศาสตร์นิวตันและทัศนศาสตร์ดั้งเดิมไม่สามารถอธิบายปรากฏการณ์ที่อัตราเร็วใกล้แสง หรือพฤติกรรมของสสารในระดับอนุภาคย่อยของอะตอมได้ การกำเนิดของ \textbf{ทฤษฎีสัมพัทธภาพพิเศษ (Special Relativity)} ของอัลเบิร์ต ไอน์สไตน์ และ \textbf{กลศาสตร์ควอนตัม (Quantum Mechanics)} ของมักซ์ พลังค์, นีลส์ บอร์, แวร์เนอร์ ไฮเซนเบิร์ก และแอร์วิน ชเรอดิงเงอร์ ได้เปิดประตูสู่ความเข้าใจเอกภพในมิติใหม่ที่ลึกซึ้งและน่าอัศจรรย์ใจ

อย่างไรก็ตาม ความท้าทายสำคัญที่สุดในการเรียนการสอนวิชาฟิสิกส์แผนใหม่ คือ \textbf{ข้อจำกัดด้านการรับรู้ทางประสาทสัมผัสของมนุษย์} เนื่องจากมนุษย์ดำเนินชีวิตในโลกสามมิติที่มีความเร็วต่ำมากเมื่อเทียบกับความเร็วแสง ทำให้มโนทัศน์อย่างการยืดของเวลา (Time Dilation), การหดสั้นของความยาว (Length Contraction), การซ้อนทับเชิงควอนตัม (Quantum Superposition), และการทะลุผ่านกำแพงศักย์ (Quantum Tunneling) กลายเป็นสิ่งที่ขัดแย้งกับสัญชาตญาณดั้งเดิม

ตำราวิชาการเรื่อง \textbf{ฟิสิกส์แผนใหม่เชิงคำนวณและการจำลองเสมือนจริง} เล่มนี้ จึงได้รับการออกแบบและสร้างสรรค์ขึ้นเพื่อแก้ปัญหาดังกล่าวอย่างบูรณาการ โดยนำเอาเทคโนโลยีความจริงเสริมและความจริงขยาย (Augmented and Extended Reality: AR/XR) บนมาตรฐานเว็บเปิด (WebXR Device API และ Three.js) มาผสานเข้ากับพลังการคำนวณทางวิทยาศาสตร์ของภาษา Python (NumPy, SciPy, Matplotlib) ทำให้นักศึกษาไม่เพียงแต่สามารถพิสูจน์สมการเชิงทฤษฎีได้อย่างแม่นยำ แต่ยังสามารถ "มองเห็น", "มีปฏิสัมพันธ์", และ "ทดลองในห้องปฏิบัติการเสมือนจริง 3 มิติ" ได้อย่างไร้ขอบเขต

เนื้อหาในตำราเล่มนี้เปิดด้วยบทนำว่าด้วยสถาปัตยกรรมการจำลองเสมือนจริง AR/XR และการคำนวณด้วย Python พร้อมทั้งเนื้อหาหลัก 6 บทเรียนสำคัญ ได้แก่
\begin{enumerate}
    \item ทฤษฎีสัมพัทธภาพพิเศษและเรขาคณิตกาล-อวกาศ 4 มิติ
    \item กลศาสตร์ควอนตัม ทวิภาวะคลื่น-อนุภาค และแบบจำลองอะตอม
    \item ฟิสิกส์ของอนุภาคและแบบจำลองมาตรฐาน (The Standard Model)
    \item ฟิสิกส์นิวเคลียร์ กัมมันตภาพรังสี และพลังงานนิวเคลียร์
    \item วิทยาการคำนวณเชิงควอนตัม คิวบิต และอัลกอริทึมควอนตัม
    \item ดาราศาสตร์ฟิสิกส์ หลุมดำ และจักรวาลวิทยาสมัยใหม่
\end{enumerate}

ผู้เขียนหวังเป็นอย่างยิ่งว่า ตำราเล่มนี้จะเป็นสะพานเชื่อมมโนทัศน์ทางฟิสิกส์ชั้นสูงเข้ากับทักษะการคำนวณเชิงตัวเลขและเทคโนโลยีดิจิทัลแห่งอนาคต สำหรับนักศึกษา ครู-อาจารย์ และนักวิจัยทุกท่าน

\vspace{1.5cm}
\begin{flushright}
    \textbf{ผู้ช่วยศาสตราจารย์ ดร.ชีวะ ทัศนา}\\
    สาขาวิชาฟิสิกส์ คณะวิทยาศาสตร์และเทคโนโลยี\\
    มหาวิทยาลัยราชภัฏรำไพพรรณี\\
    จังหวัดจันทบุรี
\end{flushright}
